\section{Discussion and Limitations}\label{sec:discussion}

\paragraph{Where the gains come from}
The experiments support a narrower claim than "keeping the network alive". Connectivity in itself
is not the bottleneck in the scenarios we studied: keeping the swarm connected at all times
(SENTINEL-strict, B3, B4) limits the search and lowers on-time reporting, and gossip between UAVs
that meet during the search already protects most information from failures. What raised timely
reporting was deciding \emph{when} a UAV should stop searching to deliver what it found, and
combining views so that occlusion does not hide confirmed victims; these are the components with
significant effects. Given the same fusion rule, the baselines remain significantly behind in the
urban 400\,m scenarios, but in a dense 50-UAV swarm they match SENTINEL and one exceeds it, so there
the fusion rule, not the mobility, is what matters. Part
of its margin in the urban default is not explained by these components
(Section~\ref{sec:eval:ablation}). The link-preservation layer (Proposition~\ref{prop:link}) did not
change the on-time share measurably, but it reduced preserved-link violations 37-fold; its role is to
make deliveries reliable when they happen.

\paragraph{Choice of $T_h$}
$T_h$ trades on-time reporting against total discovery. The value tuned on the report value $J$
(116\,s) is longer than the 60\,s deadline. It still improves on-time reporting over never
delivering (urban, $N=8$: 0.577 at $T_h=120$\,s vs.\ 0.523), but a hold time of 15\,s was better in
both channels with eight UAVs (0.599 urban, 0.501 severe; Fig.~\ref{fig:thold}). In practice $T_h$ should follow the responders' reporting deadline; adapting it online to the
observed contact rate is a natural extension.

\paragraph{Tuning sensitivity}
An earlier configuration of SENTINEL that also tethered UAVs softly while exploring helped dense
swarms but reduced coverage in the severe channel; the evaluated design does not use this tether.
Results of that configuration are released with the others.

\paragraph{Limitations}
(i)~All results are from simulation: a fast multi-agent simulator and PX4 software-in-the-loop in
Gazebo; no hardware flights were performed. (ii)~The channel model is parameterised from published
ranges and has not been calibrated against measurements in rubble. (iii)~Detection accuracy is
reported from the validation split recorded by the training runs; the detector transfers poorly to
rendered thermal scenes, which is why the in-the-loop comparison uses the detector model; that
model's parameters were chosen to reproduce the detector's qualitative behaviour and swept, not
fitted to detector outputs.
(iv)~Victims closer than the association gate (about 5\,m) can be merged into one report.
(v)~MAPPO was trained to a plateau (600 iterations) with one observation and reward design; other
MARL formulations may perform better, in particular in dense swarms, where it is already the strongest
baseline. (vi)~The exposure bound
(Proposition~\ref{prop:exposure}) could not be tested empirically because losses were rare at the
failure rates studied. (vii)~Latency was measured on a desktop GPU, not on embedded hardware.
